\chapter{赛题描述}
\section{题目与任务}
本项目对应《全国大学生嵌入式芯片与系统设计竞赛2026选题指南（中科亿海微）》中的选题一“基于以太网MAC总线实现AFDX”。赛题以以太网MAC控制器为基础，要求在FPGA内组织AFDX端系统的数据处理，包括UDP/IP协议处理、缓冲区管理、发送调度、接收管理和媒体接入控制，使应用数据经过可配置的虚拟链路传送到对端。

赛题列出的主要功能包括ARINC 664 Part 7通信、高完整性数据通信、SNMP网络管理、615A数据加卸载以及RTC时钟同步。

\section{技术参考指标}
\begin{table}[htbp]\centering\small
\caption{赛题原文中的技术参考指标}
\begin{tabularx}{\linewidth}{@{}l Y@{}}\toprule 项目 & 赛题描述\\\midrule
链路速率 & 100 Mbit/s\\
帧格式 & IEEE 802.3基础帧及ARINC 664扩展，包含VL标识与1--255序号\\
帧长 & 有效载荷通常不超过1471字节，总帧长不超过1518字节\\
时延与抖动 & 最坏端到端时延参考1--2 ms，抖动参考几十微秒\\
VL能力 & 典型端系统最多128个发送VL、4096个接收VL\\
设计接口 & 以\code{AFDX_top}为顶层，包含时钟、复位和应用字节握手接口\\\bottomrule
\end{tabularx}\end{table}

\section{报告结构}
第二章介绍端系统总体数据路径与顶层IO；第三章说明615A业务控制、控制文件编解码和文件存储；第四章说明SNMP GET、Response和Trap的实现；第五章集中讲解TFTP报文格式、块长协商与会话实现；第六章通过实际仿真解释文件、网络帧及管理消息的交互过程。

